\section{Evaluation Framework}

To answer the main research question on the efficiency gains in both storage and analysis, an 
evaluation framework is presented that has the ability to measure the impact of using different 
versions of the DeltaParq system compared to a base version of the system as it is currently 
used.

\subsection{Measurement Metrics}

The primary metrics measured are the total storage size for storing a month worth of data, which 
was determined as the maximum self-contained unit in the system design, and the analysis time for each 
day of the month given a specific analysis method. The storage size for the full month will be able to,
once compared to the storage size of a base version, determine the storage efficiency of each version, 
while the analysis time for each day of the month will be able to determine, compared to the same base 
version, how much the analysis efficiency has increased or decreased. 

While these two metrics capture the full scope necessary for answering the research question. The research 
is held to determine if a delta-based storage system is a viable option for the OPENIntel dataset. From this 
follow two extra metrics. First the time to construct the delta-compressed version which will be 
used to determine the effect of using delta-compression on the daily task of gathering and storing all the new 
data. Secondly, the time to reconstruct the full original data will be measurement to report the effects that may 
occur when needing the full data for possible migration or other use cases.

% The client commisioning the research, the OPENIntel project, 

\subsection{Measurement Validity}

\subsection{Measurement Methodology}
